\documentclass[11pt,letterpaper]{article}

\usepackage[spanish,es-noshorthands,es-tabla]{babel}
\usepackage{fontspec}
\setmainfont{Source Serif Pro}
\usepackage{microtype}
\usepackage{geometry}
\usepackage{enumitem}
\usepackage{xcolor}
\usepackage{csquotes}
\usepackage{float}
\usepackage{booktabs}
\usepackage{tikz}
\usetikzlibrary{arrows.meta,positioning,fit,backgrounds}
\usepackage{pgfplots}
\pgfplotsset{compat=1.18}
\usepackage{xurl}
\usepackage[hidelinks]{hyperref}
\usepackage[backend=biber,style=apa,sorting=nyt,maxcitenames=2]{biblatex}
\DeclareLanguageMapping{spanish}{spanish-apa}

\geometry{margin=2.3cm}
\hbadness=10000
\setlength{\emergencystretch}{1em}
\setlength{\parindent}{0pt}
\setlength{\parskip}{0.5em}
\setlist{nosep}
\addbibresource{referencias.bib}
\AtBeginBibliography{\sloppy\small}
\setlength{\bibitemsep}{0.3em}

% Colores: morado Nubank para el desarrollo capitalizado y los hitos de
% tecnología, verde para la infraestructura tecnológica, lila para el personal
% de tecnología (la parte estimada), terracota para los hitos de negocio, gris
% para la cuadrícula y azul oscuro para ejes y rótulos.
\definecolor{cCapit}{RGB}{130,10,209}
\definecolor{cInfra}{RGB}{36,122,82}
\definecolor{cPersonal}{RGB}{196,160,232}
\definecolor{cNegocio}{RGB}{186,86,30}
\definecolor{cEje}{RGB}{26,58,102}
\definecolor{cGris}{RGB}{120,116,110}

\tikzset{
  % los rótulos de la figura no se dividen con guion
  every node/.append style={execute at begin node={\hyphenpenalty=10000}},
  hito/.style={draw=#1, fill=#1!9, thick, rounded corners=2pt,
    text width=1.78cm, align=center, inner sep=2.5pt, font=\scriptsize},
  anio/.style={font=\footnotesize\bfseries, text=cEje},
  punto/.style={circle, fill=#1, inner sep=0pt, minimum size=5pt},
  carril/.style={font=\footnotesize\bfseries, text=#1, anchor=west},
}

\pgfplotsset{
  foro/.style={
    font=\small,
    axis line style={cEje},
    tick style={cEje},
    every tick label/.append style={font=\small, text=cEje},
    label style={font=\small, text=cEje},
    /pgf/number format/use comma,
    /pgf/number format/1000 sep={\,},
  },
}

\begin{document}

% --- Portada -----------------------------------------------------------------
\begin{titlepage}
\centering
\vspace*{3.2cm}
{\large\scshape Gestión del Desarrollo Tecnológico:\\
Estrategias y Análisis de Portfolio\par}
\vspace{0.5cm}
{\normalsize Módulo 6. Foro 6.3: Cuantificar la inversión en I+D\\
de una empresa\par}
\vspace{2.6cm}
{\color{cEje}\rule{0.55\textwidth}{0.8pt}\par}
\vspace{0.7cm}
{\LARGE\bfseries Cuánto invierte Nubank en I+D\par}
\vspace{0.5cm}
{\large Una estimación a partir de su 20-F, cuando la cifra\\
no aparece en los estados financieros\par}
\vspace{0.7cm}
{\color{cEje}\rule{0.55\textwidth}{0.8pt}\par}
\vfill
{\large Carlos Luis Rojas Aragonés\par}
\vspace{0.3cm}
{\normalsize Chief Technology Officer Program\par}
\vspace{0.3cm}
{\normalsize 10 de octubre de 2026\par}
\vspace*{1.5cm}
\end{titlepage}

% --- Cuerpo ------------------------------------------------------------------
Elegí Nubank (Nu Holdings), el banco digital que ha puesto a correr a la
banca tradicional en Brasil, México y Colombia. Es relevante para mi hoja de
ruta porque su núcleo bancario, NuCore, es propio y está construido sobre la
nube \parencite{nu2026}: es un caso real de lo que una plataforma moderna
puede hacer frente a sistemas heredados.

\section*{1. De una tarjeta morada a 131 millones de clientes}

David Vélez, colombiano con MBA de Stanford, llegó a Brasil como inversor de
Sequoia. Ese proyecto no prosperó, pero le dio la idea de un neobanco
\parencite{lankenau2023}. Nubank se fundó en 2013, emitió su primera tarjeta
de crédito en 2014 y llegó a México en 2019 y a Colombia en 2020. En 2021
salió a la Bolsa de Nueva York y recaudó 2\,300 millones de dólares
\parencite{nu2026}.

A finales de 2025 sumaba 131 millones de clientes: 113 millones en Brasil
(el 62\,\% de los adultos del país), 14 millones en México y más de 4 millones
en Colombia. En abril de 2025 recibió la autorización para iniciar su conversión
en banco en México y, en enero de 2026, la aprobación condicional para abrir un banco
nacional en Estados Unidos \parencite{nu2026}. La Figura~\ref{fig:linea}
resume el recorrido.

\begin{figure}[H]
\centering
\begin{tikzpicture}[x=2.12cm]
  % Eje del tiempo
  \draw[very thick, cEje, -{Stealth[length=2.8mm]}] (-0.45,0) -- (7.5,0);

  % Años, todos bajo el eje
  \foreach \x/\a in {0/2013, 1/2014, 2/2019, 3/2020, 4/2021, 5/2024,
      6/2025, 7/2026} {
    \node[anio, anchor=north] at (\x,-0.1) {\a};
  }

  % Hitos de tecnología e I+D, sobre el eje
  \foreach \x/\t in {
      3/{Compra Plataformatec y Cognitect (Clojure, Datomic)},
      5/{Compra Hyperplane, laboratorio de IA},
      6/{Publica nuFormer, su modelo fundacional},
      7/{nuFormer decide crédito en producción}} {
    \node[punto=cCapit] at (\x,0) {};
    \draw[cCapit, thick] (\x,0) -- (\x,0.42);
    \node[hito=cCapit, anchor=south] at (\x,0.42) {\t};
  }

  % Hitos de negocio, bajo el eje
  \foreach \x/\t in {
      0/{Fundación en São Paulo},
      1/{Primera tarjeta de crédito},
      2/{Llega a México},
      3/{Llega a Colombia},
      4/{Sale a bolsa en Nueva York},
      6/{Inicia su conversión a banco en México},
      7/{Aprobación condicional en EE.\,UU.}} {
    \node[punto=cNegocio] at (\x,0) {};
    \draw[cNegocio, thick] (\x,-0.5) -- (\x,-0.62);
    \node[hito=cNegocio, anchor=north] at (\x,-0.62) {\t};
  }

  % Carriles
  \node[carril=cCapit]    at (-0.45,1.55) {Tecnología e I+D};
  \node[carril=cNegocio]  at (-0.45,-2.05) {Negocio y expansión};
\end{tikzpicture}
\caption{Hitos de Nubank. Sobre el eje, las apuestas tecnológicas; bajo el
eje, el crecimiento del negocio
\parencite{nu2026,nubank2020,braithwaite2025,nu2026ia}.}
\label{fig:linea}
\end{figure}

\section*{2. Dónde invierte: una plataforma propia}

Nubank decidió construir en lugar de comprar. Su núcleo bancario corre en
microservicios sobre Amazon Web Services, escritos en Clojure y apoyados en
un libro contable inmutable con Datomic, una base de datos de la que hoy es
dueña. En 2025 promedió más de 900 despliegues de código por día hábil
\parencite{nu2026}.

Varias de sus adquisiciones fueron, en la práctica, compras de capacidad de
I+D. En 2020 adquirió Cognitect, la empresa detrás de Clojure y Datomic,
cuando tenía unos 600 ingenieros de un total de 2\,700 empleados
\parencite{nubank2020}. En 2024 compró Hyperplane, una plataforma
estadounidense de IA, y con ella trajo investigadores de Google, Meta, Apple
y Microsoft \parencite{nu2026ia}. De ahí salió nuFormer, un modelo fundacional entrenado con datos
de transacciones \parencite{braithwaite2025} que ya decide crédito en su
segmento más grande en Brasil y redujo el ciclo de experimentación de meses a
un día \parencite{nu2026ia}. Hoy el 42\,\% de sus 10\,027 empleados trabaja
en tecnología \parencite{nu2026}.

\section*{3. Cuantificar la inversión}

Nubank no reporta una línea de I+D: su 20-F indica que ese gasto está dentro
de los gastos generales y de administración, junto con back-office y
overhead \parencite{nu2026}. Por eso lo estimo como un rango con tres
componentes del propio 20-F:

\begin{enumerate}[label=\alph*), leftmargin=1.6em, itemsep=0.25em]
  \item \textbf{Desarrollo capitalizado:} intangibles desarrollados
    internamente, como nuevas funcionalidades y productos.
  \item \textbf{Infraestructura y datos} registrados en administración, es
    decir, la nube y el software que no están ligados a transacciones de
    clientes.
  \item \textbf{Personal de tecnología:} el 75\,\% del gasto de personal de
    administración (salarios, pagos en acciones y otros costos), que es la
    proporción de personal de tecnología (4\,236) frente al administrativo
    (1\,415) en 2025.
\end{enumerate}

El rango bajo suma a y b; el alto suma los tres (Tabla~\ref{tab:id}).

\begin{table}[H]
\centering
\small
\begin{tabular}{@{}lrrrrr@{}}
\toprule
Año & a) Capitalizado & b) Infraestructura & c) Personal
    & I+D estimado & \% de ingresos \\
\midrule
2023 & 165 & 175 & 449 & 340 a 789   & 4,2 a 9,8 \\
2024 & 155 & 195 & 566 & 350 a 916   & 3,0 a 8,0 \\
2025 & 286 & 255 & 595 & 541 a 1\,135 & 3,4 a 7,2 \\
\bottomrule
\end{tabular}
\caption{Inversión estimada en I+D de Nubank, en millones de US\$. Ingresos:
8\,029 (2023), 11\,517 (2024) y 15\,775 (2025) \parencite{nu2026}.}
\label{tab:id}
\end{table}

La inversión creció un 44\,\% entre 2023 y 2025 en el rango alto
(Figura~\ref{fig:inversion}), pero los ingresos casi se duplicaron, de modo
que la intensidad bajó de 4,2--9,8\,\% a 3,4--7,2\,\%. El desarrollo
capitalizado, la parte más cercana a crear producto nuevo, casi se duplicó en
2025. Por cliente, la inversión se mantuvo entre 8 y 9 dólares al año
mientras la base crecía un 39\,\%.

\begin{figure}[H]
\centering
\begin{tikzpicture}
\begin{axis}[foro,
    width=8.6cm, height=6.4cm,
    ybar stacked, bar width=26pt,
    symbolic x coords={2023,2024,2025}, xtick=data,
    ymin=0, ymax=1400,
    ylabel={Millones de US\$},
    ymajorgrids, grid style={cGris!25},
    axis x line*=bottom, axis y line*=left,
    enlarge x limits=0.25,
    legend style={draw=none, font=\small, text=cEje,
      at={(1.04,0.5)}, anchor=west, cells={anchor=west}},
    reverse legend]
  \addplot[fill=cCapit, draw=none]
    coordinates {(2023,165.1) (2024,154.6) (2025,285.7)};
  \addplot[fill=cInfra, draw=none]
    coordinates {(2023,174.6) (2024,195.1) (2025,254.9)};
  \addplot[fill=cPersonal, draw=none]
    coordinates {(2023,448.8) (2024,566.0) (2025,594.7)};
  \legend{a) Desarrollo capitalizado, b) Infraestructura y datos,
    c) Personal de tecnología (estimado)}
  \node[font=\small\bfseries, text=cEje, anchor=south]
    at (axis cs:2023,788.5) {4,2--9,8\,\%};
  \node[font=\small\bfseries, text=cEje, anchor=south]
    at (axis cs:2024,915.7) {3,0--8,0\,\%};
  \node[font=\small\bfseries, text=cEje, anchor=south]
    at (axis cs:2025,1135.3) {3,4--7,2\,\%};
\end{axis}
\end{tikzpicture}
\caption{Inversión estimada en I+D de Nubank por componente. Sobre cada barra,
la intensidad en I+D (rango bajo y alto, como porcentaje de los ingresos)
\parencite{nu2026}.}
\label{fig:inversion}
\end{figure}

\section*{4. La I+D que no aparece en la tabla}

Vélez sostiene que un modelo de crédito no se construye sin pérdidas y que
muchas de ellas funcionan como inversión en I+D, incluso como CAPEX, porque
el modelo necesita equivocarse para aprender \parencite{lankenau2023}. Si se
aceptara esa lectura, la cifra sería mucho mayor: las pérdidas esperadas de
crédito pasaron de 2\,285 a 4\,205 millones de dólares entre 2023 y 2025
\parencite{nu2026}. No las incluyo en la estimación, porque no hay forma de
separar la parte que \enquote{enseña} de la que es simple riesgo de negocio,
pero muestran que en un banco basado en datos el aprendizaje también se paga
con capital.

\section*{5. ¿Qué obtuvo a cambio?}

\begin{itemize}[leftmargin=1.2em, itemsep=0.25em]
  \item \textbf{Escala y rentabilidad.} Entre 2023 y 2025 los clientes
    pasaron de 94 a 131 millones y la utilidad neta de 1\,031 a 2\,872
    millones de dólares, casi el triple \parencite{nu2026}. En el último
    trimestre de 2025 alcanzó un ROE del 33\,\% \parencite{nu2026resultados}.
  \item \textbf{Costo de servir.} Atender a un cliente activo le cuesta
    0,8 dólares al mes, frente a 15 dólares de ingreso mensual por cliente
    activo \parencite{nu2026resultados}. En Brasil estima que su costo de
    servir y de administración por cliente activo es un 85\,\% menor que el de
    los bancos tradicionales \parencite{nu2026}.
  \item \textbf{Velocidad de producto.} Solo en 2025 lanzó más de 100
    productos y funcionalidades \parencite{nu2026}.
\end{itemize}

\section*{6. Lo que me llevo como CTO}

Nubank muestra que la intensidad en I+D no basta para juzgar a una empresa:
puede bajar como porcentaje mientras la inversión crece, si la plataforma
escala. Lo decisivo fue invertir en la base (lenguaje, base de datos, nube)
antes de necesitarla, y comprar talento cuando construirlo era más lento. La
limitación de este ejercicio es clara: mi cifra es una estimación, no un dato
auditado, y la principal fuente de incertidumbre es qué parte del personal se
dedica realmente a desarrollar.

\printbibliography[title={Referencias}]

\end{document}
